\documentclass{article}

%%% LANGUAGE: Spanish
\usepackage[spanish,english]{babel}
\addto\shorthandsspanish{\spanishdeactivate{~<>}}

%%% ENCODING
\usepackage[utf8]{inputenc}
\usepackage[T1]{fontenc}
\usepackage{lmodern}
\usepackage{microtype}

%%% PACKAGES
\usepackage{amsmath}
\usepackage{graphicx}
\usepackage{booktabs}
\usepackage{array}
\usepackage{tabularx}
\usepackage{multirow}
\usepackage{paralist}
\usepackage{xcolor}
\usepackage{tcolorbox}
\usepackage{float}
\usepackage{pdflscape}   % paginas apaisadas de verdad: el visor las gira solo
\usepackage{lipsum}
\usepackage[section]{placeins}
\usepackage[colorlinks=true,linkcolor=blue,urlcolor=blue]{hyperref}

%%% COLUMN TYPE: X justificada a la izquierda
\newcolumntype{L}{>{\raggedright\arraybackslash}X}

%%% GENERAL INFORMATION
\newcommand{\theCampus}{Campus Vitacura}
\newcommand{\theAssignment}{Sistemas de Información para la Gestión}
\newcommand{\theCourseCode}{ICN-292}
\newcommand{\theReportTitle}{Entrega 1 - UNTER}
\newcommand{\theGroup}{Grupo 2}
\newcommand{\thePar}{100}
\newcommand{\theAuthorA}{Juan Pablo Godoy Fajardo}
\newcommand{\theAuthorB}{Vicente Soto Martínez}
\newcommand{\theAuthorC}{Joaquin Tapia Mandiola}
\newcommand{\theAuthorD}{Claudio Vera Avila}
\newcommand{\theStudentIdA}{202360577-8}
\newcommand{\theStudentIdB}{202204651-1}
\newcommand{\theStudentIdC}{202260674-6}
\newcommand{\theStudentIdD}{202104575-9}
\newcommand{\theSemester}{Segundo Semestre 2026}

%%% DATE
\usepackage{datetime}
\newdateformat{monthyeardate}{%
  \monthname[\THEMONTH], \THEYEAR}

%%% PAGE DIMENSIONS
\usepackage[margin=25mm,inner=30mm]{geometry}

%%% HEADERS & FOOTERS
\usepackage{fancyhdr}
\setlength{\headheight}{45pt}
\setlength{\headsep}{12pt}

\pagestyle{fancy}
\fancyhf{}
\lhead{\includegraphics[scale=0.3]{Logo usmUSM.png}}
\chead{}
\rhead{}
\renewcommand{\headrulewidth}{0.5pt}
\fancyfoot[C]{\textbf{\thepage}}

%%% TABLE OF CONTENTS APPEARANCE
\makeatletter
\def\tableofcontents{%
 \newpage
 \centerline{\large\scshape Índice}
 \vspace*{0.3in}
 \@mkboth{CONTENTS}{CONTENTS}
 \@starttoc{toc}
}
\makeatother

%%% PLACEHOLDERS
\newcommand{\placeholdergrafico}[1]{%
\begin{figure}[H]
\centering
\begin{tcolorbox}[width=0.85\textwidth,colback=gray!10,colframe=gray!60,arc=2mm]
\centering
\vspace{2cm}
\textbf{[ ESPACIO PARA GRÁFICO ]}\\[0.3cm]
\textit{#1}\\[0.3cm]
\vspace{2cm}
\end{tcolorbox}
\end{figure}
}

\newcommand{\placeholdertexto}{%
Lorem ipsum dolor sit amet, consectetur adipiscing elit. Integer vitae sem non justo tempor facilisis. Vestibulum ante ipsum primis in faucibus orci luctus et ultrices posuere cubilia curae; Proin aliquet, justo vitae consequat dictum, lorem velit tincidunt neque, sed consequat nunc libero non justo.
}

%%% MARCA DE DATO PENDIENTE
%%% Cuando esten todos completos, reemplazar por: \newcommand{\pend}[1]{}
\newcommand{\pend}[1]{\textcolor{red}{\textbf{[#1]}}}

\hypersetup{pdftitle={Entrega 1 - UNTER}}

%------------------------------------------------------------------------------
%%%% TITLE PAGE
%------------------------------------------------------------------------------
\begin{document}
\selectlanguage{spanish}
\hypersetup{pageanchor=false}

\begin{titlepage}
\thispagestyle{empty}

\vspace*{1.2cm}

\begin{center}
\includegraphics[scale=0.5]{Logo usmUSM.png}
\end{center}

\vspace{2.2cm}

\begin{center}
\hrule
\vspace{2mm}
{\Huge \bfseries \theReportTitle}\\[0.5cm]
{\Large \theCourseCode}\\[0.3cm]
{\Large \scshape \theAssignment}
\vspace{5mm}
\hrule
\vspace{1cm}

%%%% AUTHORS %%%%
\centering

\textbf{\theGroup} \quad -- \quad \textbf{Paralelo \thePar} \\
\vspace{0.4cm}
\textbf{Integrantes:}\\[0.3cm]

\textbf{\theAuthorA} \\
{\small \theStudentIdA}\\[0.3cm]

\textbf{\theAuthorB} \\
{\small \theStudentIdB}\\[0.3cm]

\textbf{\theAuthorC} \\
{\small \theStudentIdC}\\[0.3cm]

\textbf{\theAuthorD} \\
{\small \theStudentIdD}

\vspace{0.5cm}
\textbf{Repositorio en GitHub:}\\[0.2cm]
{\small \url{https://github.com/JotaReyyyyyyyyyyyy/ICN292-E1-Unter}}

\vfill
\textsc{\theSemester}\\[0.5cm]
\textsc{\theCampus}
\end{center}
\end{titlepage}
\hypersetup{pageanchor=true}
\pagenumbering{arabic}

\tableofcontents
\newpage

%==============================================================
\section{Portada y ficha del equipo}

\begin{table}[H]
\centering
\small
\begin{tabular}{lllp{4.6cm}}
\toprule
\textbf{Integrante} & \textbf{Rol USM} & \textbf{RUT} & \textbf{Rol en el equipo} \\
\midrule
\theAuthorA & \theStudentIdA & 21.633.620-8 & Evidencia y gobernanza de datos \\
\theAuthorB & \theStudentIdB & 21.270.145-9 & Interfaz y KPI \\
\theAuthorC & \theStudentIdC & 21.407.361-7 & Modelo de datos y automatizacion \\
\theAuthorD & \theStudentIdD & 22.231.192-6 & Base de datos y reproducibilidad \\
\bottomrule
\end{tabular}
\end{table}

\noindent\textbf{Paralelo:} \thePar

\noindent\textbf{Repositorio GitHub:} \url{https://github.com/JotaReyyyyyyyyyyyy/ICN292-E1-Unter}

%==============================================================
\section{PYME y problema (evidencia)}

\subsection{Identificación de la PYME}

\begin{table}[H]
\centering
\begin{tabularx}{\textwidth}{lL}
\toprule
\textbf{Campo} & \textbf{Valor} \\
\midrule
Nombre & Unter SpA \\
Rubro & Distribuidora de equipos de protección personal (EPP) \\
Tamaño & Microempresa: 2 trabajadores \\
Ubicación & La Concepción 81, Providencia, Santiago \\
\bottomrule
\end{tabularx}
\end{table}

\noindent La operación completa de Unter recae en dos personas:

\begin{itemize}
\item \textbf{Ariel} --- Gerente General. Área comercial, catálogo, especificaciones técnicas, prospección y atención de clientes.
\item \textbf{Jani} --- Gerente de Finanzas. Tesorería, precios, condiciones de pago y plazos, y la totalidad del trato y la negociación con proveedores.
\end{itemize}

\subsection{Evidencia de existencia}

Se realizó una primera instancia de levantamiento en modalidad informal. La entrevista formal es un audio que se encuentra en el GitHub junto con la transcripción

\subsection{Problema de negocio}

Hoy Unter gestiona las consultas de clientes, la prospección comercial y la administración técnica de forma manual y sin un sistema que logre centralizar toda la información. Esto genera respuestas lentas, oportunidades de venta que se pierden por el camino e información poco consistente, lo que termina afectando tanto la calidad del servicio como la capacidad de la empresa para crecer.

\subsubsection*{Impacto}

\begin{itemize}
\item \textbf{Tiempo:} buena parte de las horas de Jani y Ariel se va en buscar información dispersa y en coordinarse manualmente entre ellos. Eso encarece cada operación y les deja menos tiempo para dedicarlo a la gestión estratégica del negocio.
\item \textbf{Costo:} se pierde dinero que se podría estar ganando, porque las oportunidades de venta se escapan cuando las respuestas tardan demasiado y los esfuerzos comerciales se hacen de forma reactiva y poco eficiente.
\item \textbf{Error:} al compartir la documentación original de los proveedores existe el riesgo de que el cliente termine comprándole directamente a ellos, y no a Unter. Además, sin un repositorio técnico ordenado, es muy fácil entregar especificaciones inconsistentes o desactualizadas.
\item \textbf{Servicio:} la experiencia del cliente empeora, porque la atención está dividida entre varias personas, los tiempos de espera son impredecibles y no hay forma de hacer seguimiento a todo el ciclo de consulta y venta de principio a fin.
\end{itemize}

\subsubsection*{Indicadores del problema}

Los siguientes indicadores le ponen números al problema que describimos y sirven como punto de partida sobre el cual mediremos los KPI en la Entrega 2.

\begin{table}[H]
\centering
\begin{tabularx}{\textwidth}{Lcc}
\toprule
\textbf{Indicador} & \textbf{Valor actual} & \textbf{Fuente} \\
\midrule
Tiempo de búsqueda de datos de clientes & 1.5 h & Entrevista \\
Consultas de clientes recibidas al mes & 12 & Entrevista \\
Tiempo promedio de respuesta a una consulta & 30 min & Entrevista \\
Peor tiempo de respuesta registrado & No responder & Entrevista \\
Consultas al mes no resueltas por falta de información & 0 & Entrevista \\
Productos distintos en catálogo & 3000 & Entrevista \\
Productos con ficha técnica propia de Unter & 0 & Entrevista \\
Tiempo empleado en elaborar una ficha & 20 min & Entrevista \\
\bottomrule
\end{tabularx}
\end{table}

\subsection{Objetivo del SIG propuesto}

Implementar un sistema de información que automatice la atención de consultas de clientes, la prospección comercial y la normalización de la documentación técnica, de manera que responderle a un cliente ya no dependa de si una persona en particular está disponible o se acuerda de algo, sino de un dato registrado al que se le puede hacer seguimiento.

\medskip
\noindent\textbf{Meta de desempeño:} reducir en un 30\,\% el tiempo promedio de respuesta y la carga de trabajo que conlleva, de manera que la operación pueda crecer sin necesidad de sumar más personas a las dos que hay hoy. Esta meta se verificará en la Entrega 2 comparándola con la línea base que levantamos en la tabla de indicadores anterior.

%==============================================================
\section{Requerimientos desde cero}

\subsection{Actores y roles}

\subsubsection*{Actores humanos}

\begin{table}[H]
\centering
\begin{tabularx}{\textwidth}{llL}
\toprule
\textbf{Actor} & \textbf{Tipo} & \textbf{Rol} \\
\midrule
Ariel & Interno & Gerente General. Área comercial, catálogo, especificaciones técnicas, prospección y atención. Asume cualquier rol adicional que se identifique. \\
Jani & Interno & Gerente de Finanzas. Tesorería, precios, condiciones de pago y plazos, y trato y negociación con proveedores. \\
Cliente & Externo & Empresa o persona que consulta por productos, precios o especificaciones técnicas. \\
Proveedor & Externo & Empresa que abastece los productos y origina la documentación técnica. \\
\bottomrule
\end{tabularx}
\end{table}

\subsubsection*{Sistemas}

\begin{table}[H]
\centering
\begin{tabularx}{\textwidth}{lL}
\toprule
\textbf{Sistema} & \textbf{Función} \\
\midrule
Asistente de IA (bot) & Atiende el chat, interpreta la consulta, responde desde la base de información aprobada y deriva cuando no puede resolver. \\
Motor de búsqueda de clientes (N8N) & Recorre fuentes abiertas, recopila y organiza empresas y contactos potenciales. \\
Generador de fichas técnicas & Extrae y normaliza atributos desde la documentación del proveedor y arma borradores con identidad de Unter. \\
\bottomrule
\end{tabularx}
\end{table}

\subsection{Alcance}

\begin{table}[H]
\centering
\small
\begin{tabularx}{\textwidth}{lLL}
\toprule
\textbf{Módulo} & \textbf{Dentro del alcance (in)} & \textbf{Fuera del alcance (out)} \\
\midrule
Asistente de IA (Bot) & Registro e historial de consultas, respuestas y derivaciones. & Resolución autónoma de negociaciones, precios o compromisos de pago y entrega. \\
\addlinespace
Asistente de IA (Bot) & Derivación a Jani o Ariel según el tema. & Incorporación de un tercer responsable o nivel de escalamiento. \\
\addlinespace
Motor de búsqueda de clientes (N8N) & Búsqueda y organización de empresas y contactos potenciales. & Envío automático de campañas o mensajes: Ariel realiza los envíos manualmente. \\
\addlinespace
Generador de fichas técnicas & Generación de borradores de fichas técnicas con identidad de Unter a partir de documentación del proveedor. & Fabricación de productos, modificación de sus especificaciones o presentación de Unter como fabricante. \\
\bottomrule
\end{tabularx}
\end{table}

\subsection{Requisitos funcionales}

Los requisitos se presentan ordenados por prioridad descendente, donde 1 corresponde a la mayor prioridad. La columna Razón muestra cómo se conecta cada requisito con el problema que levantamos.

\begin{table}[H]
\centering
\small
\begin{tabularx}{\textwidth}{ccLL}
\toprule
\textbf{Prio.} & \textbf{ID} & \textbf{Requisito} & \textbf{Razón} \\
\midrule
1 & RF01 & Recibir consultas mediante chat y registrar cada consulta. & Es la entrada principal del sistema y la base del módulo de atención. \\
\addlinespace
2 & RF02 & Responder consultas únicamente con información y fichas técnicas aprobadas. & Es la función central del asistente y evita respuestas basadas en información no validada. \\
\addlinespace
3 & RF03 & Registrar información faltante y derivar la consulta cuando no pueda responderse. & Evita que el sistema invente información y asegura continuidad de atención. \\
\addlinespace
4 & RF04 & Derivar consultas a Jani o Ariel según el tipo de asunto. & Automatiza una regla de negocio concreta y reduce trabajo de coordinación. \\
\addlinespace
5 & RF05 & Buscar potenciales clientes según criterios definidos por Ariel. & Representa la funcionalidad principal del módulo de prospección. \\
\addlinespace
6 & RF06 & Guardar datos de empresas y contactos junto con fuente y fecha de obtención. & Permite construir una base comercial útil y trazable. \\
\addlinespace
7 & RF07 & Generar un borrador de ficha técnica Unter desde documentación del proveedor. & Es la funcionalidad de mayor valor del módulo de fichas técnicas. \\
\addlinespace
8 & RF08 & Permitir que Ariel revise, corrija y apruebe las fichas antes de utilizarlas. & Garantiza control humano y conecta directamente las fichas con el asistente. \\
\bottomrule
\end{tabularx}
\end{table}

\subsection{Requisitos no funcionales}

\begin{table}[H]
\centering
\small
\begin{tabularx}{\textwidth}{ccLL}
\toprule
\textbf{Prio.} & \textbf{ID} & \textbf{Requisito} & \textbf{Razón} \\
\midrule
1 & RNF01 & El asistente no debe entregar nunca una especificación técnica o un precio que no esté en la base: ante la ausencia del dato debe declararlo y derivar. & Un dato inventado entregado con la marca de Unter compromete la responsabilidad de la empresa frente al cliente. \\
\addlinespace
2 & RNF02 & Ninguna ficha técnica puede publicarse sin validación humana previa. & Desde que la ficha lleva la marca de Unter, el error deja de ser atribuible al proveedor. \\
\addlinespace
3 & RNF03 & Todo dato de contacto recopilado debe conservar su fuente y fecha de obtención, y debe existir un mecanismo de baja a solicitud del titular. & Cumplimiento de la Ley 21.719 sobre protección de datos personales. \\
\addlinespace
4 & RNF04 & El sistema debe ser operable por una persona sin formación técnica. & Unter no cuenta con personal de TI. \\
\addlinespace
5 & RNF05 & El registro de una respuesta o de un contacto enviado debe requerir un solo paso. & Con una dotación de dos personas, todo registro que exija más de un paso no se realiza, y sin registro el sistema pierde su valor. \\
\bottomrule
\end{tabularx}
\end{table}

%==============================================================
\section{Procesos BPMN}

Modelamos los tres procesos críticos del caso, cada uno en su versión as-is y to-be. Los participantes son las dos personas de Unter y el sistema propuesto. La regla para repartir el trabajo es simple: todo lo que tenga que ver con dinero, condiciones comerciales o negociación con proveedores queda en finanzas y proveedores, y el resto queda en operaciones. Por eso la generación de fichas técnicas, que es un trámite de documentos y no una negociación, ocurre completamente en operaciones. Como es una empresa de solo dos personas, no hay un tercer nivel al cual escalar.

En los diagramas as-is el carril del sistema no existe, y esa ausencia es el diagnóstico.

\subsection{Atención de consultas de clientes}

Hoy la consulta entra por el canal que el cliente tenga a mano y se resuelve según quién la reciba y qué alcance a recordar. Cuando falta el dato, hay que preguntarle al proveedor y el cliente queda esperando sin un plazo claro. Como nada queda registrado, esa misma consulta se termina resolviendo de nuevo desde cero la próxima vez.

En el to-be el chat pasa a ser la única entrada: el asistente lee la consulta, revisa la base y responde cuando tiene la información suficiente. Si no la tiene, junta lo que hay disponible, genera el aviso y le pasa el caso al gestor comercial con todo el contexto ya armado.

Los diagramas as-is y to-be de este proceso están en las figuras~\ref{fig:aten-asis} y~\ref{fig:aten-tobe} del Anexo~\ref{anexo:bpmn}.

\subsection{Búsqueda de clientes}

La prospección hoy es manual y sin memoria. Se busca a mano, se escribe cada mensaje desde cero y el resultado termina anotado en una planilla o simplemente en ninguna parte. Tampoco queda registro de a quién ya se contactó.

En el to-be esa búsqueda se transforma en una gestión ordenada con criterios definidos. El sistema recopila y guarda los datos de contacto junto con su origen, califica a los prospectos y le entrega la lista al gestor comercial. El envío del primer contacto sigue siendo manual, ya que el sistema prepara todo, pero es la persona quien decide y envía.

Los diagramas as-is y to-be de este proceso están en las figuras~\ref{fig:busq-asis} y~\ref{fig:busq-tobe} del Anexo~\ref{anexo:bpmn}.

\subsection{Fichas técnicas}

La generación de fichas es un proceso interno, no lo inicia el cliente. Hoy tiene dos puntos de partida que están desconectados entre sí. Cuando entra un producto nuevo al inventario, el documento del proveedor se guarda en un correo o una carpeta y ahí queda, así que el producto entra al catálogo sin una ficha propia. Y cuando después se necesita esa ficha, alguien tiene que buscar el documento, pedírselo al proveedor si no lo encuentra, y armarla a mano en PowerPoint. Esa ficha se usa solo para la consulta puntual que la motivó y después se descarta, porque no queda ni en el catálogo ni en ningún repositorio. El resultado es que todo el trabajo se repite de nuevo la próxima vez que alguien pregunta por el mismo producto.

El to-be transforma eso en un creador de fichas semi-automático que trabaja sobre el inventario y no sobre las consultas. Tiene dos formas de activarse, una automática, cuando se agrega un producto nuevo, y otra a pedido, cuando operaciones lanza una revisión periódica para generar de una vez todas las fichas que falten. A partir de ahí el sistema detecta los productos que no tienen ficha Unter, busca la documentación del proveedor, extrae y ordena los atributos, y genera la ficha. La aprobación siempre queda en manos de una persona, ya que ninguna ficha se publica sin que operaciones la revise primero. Si falta documentación, operaciones se la pide al proveedor y el producto queda marcado como sin respaldo. El ciclo se repite hasta que no queden productos pendientes, de modo que el catálogo tiende a estar siempre con sus fichas al día.
Los diagramas as-is y to-be de este proceso están en las figuras~\ref{fig:fich-asis} y~\ref{fig:fich-tobe} del Anexo~\ref{anexo:bpmn}.

\subsection{Mejoras que introduce el SIG}

\textbf{Qué se automatiza.} La recepción y clasificación de las consultas, la búsqueda de la respuesta en la base, la recopilación de prospectos y la eliminación de los repetidos, y la extracción y ordenamiento de los atributos técnicos del documento del proveedor. Son todas tareas que hoy les consumen tiempo a las dos personas sin que aporten criterio real.

\textbf{Qué se controla.} El plazo de las derivaciones, el seguimiento de los prospectos que no han respondido, y qué productos del catálogo se quedaron sin ficha propia. En el as-is nada de esto está controlado, porque todo depende de que alguien se acuerde.

\textbf{Qué se mide.} Las consultas que recibe y resuelve el asistente, el tiempo de respuesta promedio y el del peor caso, los prospectos contactados y los que se convirtieron, y cuánto del catálogo alcanza a tener ficha propia. Ninguno de estos indicadores existe hoy, y son justamente la línea base de los KPI de la Entrega 2.

\textbf{Qué no cambia.} Las decisiones las siguen tomando las personas. El sistema prepara, prioriza y registra, pero aprobar una ficha y enviar un primer contacto siempre queda en manos de alguien.

%==============================================================
\section{Modelo de datos preliminar (ER)}

El modelo sostiene los tres procesos to-be con nueve entidades. Se usa la notación simplificada de la asignatura: cajas con atributos y tipo, clave primaria marcada con \texttt{\#} y \texttt{(ip)}, foráneas con \texttt{(fk)}, y la relación N:M resuelta con una entidad puente.

\begin{landscape}
\thispagestyle{plain}   % el encabezado con el logo, girado, queda de canto
\begin{figure}[H]
\centering
\includegraphics[width=\linewidth,height=0.9\textheight,keepaspectratio]{er_unter.png}
\caption{Modelo entidad--relación preliminar de Unter.}
\end{figure}
\end{landscape}

\subsection{Entidades}

\begin{table}[H]
\centering
\small
\begin{tabularx}{\textwidth}{llL}
\toprule
\textbf{Entidad} & \textbf{PK candidata} & \textbf{Qué representa} \\
\midrule
CLIENTE & ID\_CLIENTE & Empresa externa, sea prospecto o cliente. Se modela una sola entidad porque un prospecto que compra sigue siendo la misma empresa. \\
CONTACTO & ID\_CONTACTO & Persona dentro de una empresa. Separada porque el derecho de baja de la Ley 21.719 es de la persona, no de la organización. \\
CONSULTA & ID\_CONSULTA & Cada interacción de un cliente con el asistente, con su canal, categoría y tiempos. \\
DERIVACION & ID\_DERIVACION & Traspaso de una consulta a una persona, con asignado, plazo y estado. \\
PRODUCTO & SKU & Ítem del catálogo. El SKU es único y esa unicidad es lo que detecta la carga duplicada. \\
FICHA\_TECNICA & ID\_FICHA & Documentación con marca Unter. Entidad propia y versionada, no un campo de PRODUCTO, porque la aprobación y el versionado son parte del modelo. \\
PROVEEDOR & ID\_PROV & Quien abastece los productos y origina la documentación técnica. \\
CAMPANA & ID\_CAMPANA & Búsqueda de prospectos con criterios definidos. \\
PROSPECCION & ID\_CAMPANA + ID\_CLIENTE & Entidad puente: resuelve el N:M entre campañas y empresas, y guarda lo que ocurrió en esa campaña. \\
\bottomrule
\end{tabularx}
\end{table}

\subsection{Relaciones y cardinalidades}

\begin{table}[H]
\centering
\small
\begin{tabularx}{\textwidth}{lcL}
\toprule
\textbf{Relación} & \textbf{Card.} & \textbf{Lectura} \\
\midrule
CLIENTE -- CONTACTO & 1:N & Una empresa tiene varios contactos; cada contacto pertenece a una empresa. \\
CONTACTO -- CONSULTA & 1:N & Un contacto puede hacer muchas consultas. \\
CONSULTA -- DERIVACION & 1:N & Una consulta puede derivarse más de una vez. \\
PRODUCTO -- CONSULTA & 1:N & Un producto puede ser objeto de muchas consultas; una consulta apunta a lo más a un producto. \\
PRODUCTO -- FICHA\_TECNICA & 1:N & Un producto acumula versiones de ficha, con una sola vigente. \\
PROVEEDOR -- PRODUCTO & 1:N & Un proveedor abastece varios productos. \\
CAMPANA -- PROSPECCION & 1:N & Una campaña evalúa muchas empresas. \\
CAMPANA -- CLIENTE & N:M & Resuelta mediante la entidad puente PROSPECCION. \\
\bottomrule
\end{tabularx}
\end{table}

\subsection{Trazabilidad proceso $\leftrightarrow$ datos}

\begin{table}[H]
\centering
\small
\begin{tabularx}{\textwidth}{LL}
\toprule
\textbf{Paso del proceso to-be} & \textbf{Qué lo sostiene en el modelo} \\
\midrule
Reconocer a un cliente ya registrado. & CLIENTE + CONTACTO. \\
Responder desde la base sin intervención humana. & FICHA\_TECNICA vigente asociada al PRODUCTO. \\
Derivar con contexto y plazo controlado. & DERIVACION, con \texttt{ASIGNADO\_A}, \texttt{PLAZO} y \texttt{ESTADO}. \\
Medir el tiempo de respuesta. & CONSULTA: \texttt{FECHA\_INGRESO} y \texttt{FECHA\_RESPUESTA}. \\
Detectar productos sin ficha propia. & PRODUCTO.\texttt{ESTADO\_DOC}. \\
Guardar el origen de cada dato de contacto. & CONTACTO: \texttt{FUENTE\_DATO} y \texttt{FECHA\_OBTENCION}. \\
Bloquear a quien pide no ser contactado. & CLIENTE.\texttt{NO\_CONTACTAR}. \\
No volver a prospectar una empresa ya trabajada. & PROSPECCION histórica por campaña. \\
Publicar una ficha solo tras aprobación. & FICHA\_TECNICA: \texttt{ESTADO} y \texttt{VIGENTE}. \\
\bottomrule
\end{tabularx}
\end{table}

El punto de cierre entre los tres procesos está en PRODUCTO. El proceso de fichas recorre el inventario y llena FICHA\_TECNICA; el de atención lee la ficha vigente para responder sin intervención humana; y el de prospección alimenta CLIENTE y CONTACTO, que son quienes después hacen las consultas. Esa dependencia es la razón por la que los tres módulos son un solo sistema y no tres proyectos separados.

%==============================================================
\section{Arquitectura lógica y stack}

\subsection{Componentes}

\begin{table}[H]
\centering
\small
\begin{tabularx}{\textwidth}{lLL}
\toprule
\textbf{Capa} & \textbf{Componente} & \textbf{Función} \\
\midrule
Captura de datos & Chat del cliente, panel interno y carga de documentos del proveedor. & Punto único de entrada de consultas, definición de campañas y alta de documentación técnica. \\
\addlinespace
Automatización & Motor de flujos (n8n u homólogo). & Ejecuta la prospección y la generación de fichas conectadas al BPMN to-be. \\
\addlinespace
Procesamiento & Reconocimiento de producto, extracción y normalización de atributos. & Convierte el documento del proveedor en atributos con formato Unter. \\
\addlinespace
Persistencia & Base de datos relacional con el modelo de la sección anterior. & Fuente única de productos, fichas, consultas y prospectos. \\
\addlinespace
Analítica y KPI & Tablero de indicadores. & Tiempo de respuesta, consultas resueltas por el asistente, cobertura de fichas y conversión de prospectos. \\
\addlinespace
Interfaz & Navegador en localhost. & Chat para el cliente y panel para las dos personas de Unter. \\
\bottomrule
\end{tabularx}
\end{table}

\subsection{Stack tentativo para la Entrega 2}

\begin{table}[H]
\centering
\small
\begin{tabularx}{\textwidth}{lL}
\toprule
\textbf{Componente} & \textbf{Elección tentativa} \\
\midrule
Lenguaje & Python 3.11, con dependencias fijadas en \texttt{requirements.txt}. \\
Interfaz y API & Streamlit para el panel y el tablero; FastAPI si el chat requiere un servicio aparte. \\
Base de datos & PostgreSQL, con SQLite como alternativa si se prioriza levantar sin contenedores. \\
Automatización & n8n, ejecutado en local. \\
Reproducibilidad & \texttt{docker compose up} más \texttt{schema.sql} y \texttt{seed.sql} versionados en el repositorio. \\
\bottomrule
\end{tabularx}
\end{table}

El criterio de selección no es la sofisticación sino que un tercero pueda levantar la solución siguiendo solo el \texttt{README}: clonar, instalar dependencias fijadas, cargar el esquema y los datos, levantar la interfaz y ver los KPI. Los datos de carga serán sintéticos o anonimizados, nunca los reales de Unter.

\section{Plan de Entrega 2 y riesgos}

\subsection{Marco normativo y puesta en marcha}

\begin{table}[H]
\centering
\small
\renewcommand{\arraystretch}{1.2}
\begin{tabularx}{\textwidth}{@{}p{0.24\textwidth}L@{}}
\toprule
\textbf{Aspecto} & \textbf{Descripción} \\
\midrule
Entrada en vigencia & La Ley 21.719 entra en vigencia el \textbf{1 de diciembre de 2026}, diez días después del cierre de la Entrega 2, por lo que el sistema debe cumplirla desde su puesta en marcha. \\
\addlinespace
Régimen para PYMEs & Unter, como PYME, queda cubierta por el artículo sexto transitorio: durante el primer año la Agencia puede amonestar en lugar de multar, pero las obligaciones son las mismas. \\
\bottomrule
\end{tabularx}
\end{table}

\subsection{Hitos y responsables}

\begin{table}[H]
\centering
\small
\renewcommand{\arraystretch}{1.2}
\begin{tabularx}{\textwidth}{@{}>{\raggedright\arraybackslash}p{0.19\textwidth}L>{\raggedright\arraybackslash}p{0.20\textwidth}@{}}
\toprule
\textbf{Plazo} & \textbf{Hito} & \textbf{Responsable} \\
\midrule
12--30 sep & Entrevista formal con acta e indicadores del problema; modelo de datos en 3FN y \texttt{schema.sql}. & Godoy / Tapia \\
\addlinespace
1--14 oct & \texttt{seed.sql} con datos anonimizados y tres consultas de negocio. & Vera \\
\addlinespace
15--28 oct & Interfaz en localhost y tablero de KPI. & Soto \\
\addlinespace
29 oct--7 nov & Automatización (n8n) conectada al BPMN to-be y \path{docs/07-etica-ley-21719.md}. & Tapia / Godoy \\
\addlinespace
8--19 nov & \texttt{README} reproducible, informe actualizado y \texttt{CHANGELOG.md}. & Vera / equipo \\
\addlinespace
21 nov & Entrega en Aula y GitHub, 12:30. & Todo el equipo \\
\bottomrule
\end{tabularx}
\end{table}

\subsection{Riesgos de datos personales en la operación del sistema}

\begin{table}[H]
\centering
\small
\renewcommand{\arraystretch}{1.2}
\begin{tabularx}{\textwidth}{@{}LL@{}}
\toprule
\textbf{Riesgo} & \textbf{Mitigación} \\
\midrule
El módulo de prospección recopila datos de contacto sin una base de licitud que lo justifique. & Interés legítimo como base, con test de balanceo documentado y registro de actividades de tratamiento. \\
\addlinespace
Un contacto pide no ser contactado y el sistema vuelve a incluirlo en la siguiente campaña. & RF-21: marca \texttt{no\_contactar} en la empresa y bloqueo automático en toda campaña futura. \\
\addlinespace
La base acumula prospectos descartados de forma indefinida, sin finalidad vigente. & Plazo de conservación definido y purga periódica de los prospectos sin actividad. \\
\addlinespace
El asistente entrega al cliente información de otro cliente, o especificaciones que no existen en la base. & Identificación del interlocutor antes de responder, y RNF-01: solo responde desde fichas aprobadas. \\
\addlinespace
Acceso no autorizado a la base de contactos y fichas. & Control de acceso por usuario y procedimiento de notificación a la Agencia dentro de 72 horas. \\
\bottomrule
\end{tabularx}
\end{table}

\appendix
\section{Diagramas BPMN\label{anexo:bpmn}}

\begin{landscape}
\thispagestyle{plain}
\begin{figure}[H]
\centering
\includegraphics[width=\linewidth,height=0.9\textheight,keepaspectratio]{atencion-ASIS.png}
\caption{Atención de consultas de clientes --- as-is.\label{fig:aten-asis}}
\end{figure}
\end{landscape}

\begin{landscape}
\thispagestyle{plain}
\begin{figure}[H]
\centering
\includegraphics[width=\linewidth,height=0.9\textheight,keepaspectratio]{atencion-TOBE.png}
\caption{Atención de consultas de clientes --- to-be, con el asistente y la base de datos.\label{fig:aten-tobe}}
\end{figure}
\end{landscape}

\begin{landscape}
\thispagestyle{plain}
\begin{figure}[H]
\centering
\includegraphics[width=\linewidth,height=0.9\textheight,keepaspectratio]{busqueda-ASIS.png}
\caption{Búsqueda de clientes --- as-is.\label{fig:busq-asis}}
\end{figure}
\end{landscape}

\begin{landscape}
\thispagestyle{plain}
\begin{figure}[H]
\centering
\includegraphics[width=\linewidth,height=0.9\textheight,keepaspectratio]{busqueda-TOBE.png}
\caption{Búsqueda de clientes --- to-be, con campañas y base de prospectos.\label{fig:busq-tobe}}
\end{figure}
\end{landscape}

\begin{landscape}
\thispagestyle{plain}
\begin{figure}[H]
\centering
\includegraphics[width=\linewidth,height=0.9\textheight,keepaspectratio]{fichas-ASIS.png}
\caption{Generación de fichas técnicas --- as-is.\label{fig:fich-asis}}
\end{figure}
\end{landscape}

\begin{landscape}
\thispagestyle{plain}
\begin{figure}[H]
\centering
\includegraphics[width=\linewidth,height=0.9\textheight,keepaspectratio]{fichas-TOBE.png}
\caption{Generación de fichas técnicas --- to-be, creador semi-automático que recorre el inventario.\label{fig:fich-tobe}}
\end{figure}
\end{landscape}

\clearpage

\section{Anexos (capturas, acta)\label{anexo:material}}
\textit{Por completar: capturas de respaldo y acta de reunión.}

\end{document}
